\documentclass[tikz,border=10pt]{standalone}
\usepackage{ctex}
\usepackage{amsmath,amssymb}
\usetikzlibrary{positioning,arrows.meta,shapes.geometric,fit}

\begin{document}
\begin{tikzpicture}[
  font=\small,
  box/.style={draw, rounded corners=2pt, align=center, minimum height=1.2cm, minimum width=2.8cm, inner sep=4pt},
  pre/.style={box, fill=blue!8},
  mid/.style={box, fill=orange!12},
  outbox/.style={box, fill=green!10},
  blackbox/.style={draw, dashed, thick, rounded corners=5pt, inner sep=8pt},
  arrow/.style={-{Stealth[length=3mm]}, thick},
  note/.style={font=\scriptsize\bfseries, text=black!70}
]

% ========== 输入 / 预处理 ==========
\node[pre] (q) {L 个快拍\\ 1-bit 复信号\\ \(q_1,\ q_2\)};
\node[pre, right=1.0cm of q] (tau) {已知均匀抖动\\ \(\tau_1,\ \tau_2\)\\ \(T=3\)};
\node[pre, right=1.0cm of tau] (cov) {1-bit 协方差\\ \(\hat{R}\)\\ 上三角 272 维};
\draw[arrow] (q) -- (tau);
\draw[arrow] (tau) -- (cov);
\node[note, above=0.05cm of q] {输入 / 预处理（线性空间）};

% ========== 网络（黑盒） ==========
\node[mid, right=1.4cm of cov] (in) {in\_layer\\ \(272 \rightarrow 64\)};
\node[mid, below=0.45cm of in] (reshape) {reshape\\ \((2, 32)\)};
\node[mid, below=0.45cm of reshape] (conv) {6\(\times\) Conv1d\\ + BN + ReLU};
\node[mid, below=0.45cm of conv] (attn) {CrossAttention\\ 4 heads};
\node[mid, below=0.45cm of attn] (out) {out\_layer\\ \(64 \rightarrow 32\)};

\draw[arrow] (cov) -- (in);
\draw[arrow] (in) -- (reshape);
\draw[arrow] (reshape) -- (conv);
\draw[arrow] (conv) -- (attn);
\draw[arrow] (attn) -- (out);

\node[blackbox, fit=(in)(out)] (bb) {};
\node[note, above=0.05cm of bb] {网络（核心黑盒：可学习非线性映射）};

% ========== 输出 / DOA ==========
\node[outbox, right=1.4cm of out] (h) {\(h \in \mathbb{R}^{32}\)};
\node[outbox, right=1.4cm of h] (sp) {空间谱\\ \(sp(\theta) = |a^H(\theta) h|^2\)};
\node[outbox, right=1.4cm of sp] (doa) {峰值搜索\\ DOA 估计};
\draw[arrow] (out) -- (h);
\draw[arrow] (h) -- (sp);
\draw[arrow] (sp) -- (doa);
\node[note, above=0.05cm of h] {输出（字典/谱/峰值搜索，非学习部分）};

% ========== 底部说明 ==========
\node[note, below=0.5cm of doa] {关键：1-bit 非线性输入 \(\rightarrow\) 抖动协方差 \(\rightarrow\) 线性空间 \(\rightarrow\) 网络 \(\rightarrow\) 空间谱 \(\rightarrow\) DOA};

\end{tikzpicture}
\end{document}
